\subsection{Robustness checks for the family of pre-committed games in Section~\ref{sec:ext-valid}}
\label{app:robustness_checks}

\begin{table}[h!]
\begin{center}
\caption{Summary statistics for absolute predictive accuracy of different models ($\varepsilon = 0$).} 
\label{tab:match_rates}
\footnotesize
\renewcommand{\arraystretch}{1.25}
\input{tables/abs_stats_table.tex}
\end{center}
\begin{footnotesize}
\begin{singlespace}
\vspace*{-0.05in}
\emph{Notes:} This table reports summary statistics for the absolute predictive accuracy of different models in Section~\ref{sec:ext-valid}.
These are the raw results without any smoothing.
The first row corresponds to the proportion of human subjects who chose the most likely strategy for the given model in the columns.
The second row corresponds to the proportion of human subjects who chose the top 3 most likely strategies for the given model in the columns.
The third is the proportion of human subjects who chose a strategy with a positive probability for the given model in the columns.
The fourth is the proportion of games in which any human chose a strategy with a positive probability for the given model in the columns.
The fifth is the proportion of games in which all humans chose a strategy with a positive probability for the given model in the columns.
Note that the cognitive hierarchy model, by design, provides positive probability for all strategies (even without smoothing) because the level-$0$ player chooses uniformly at random.
None of the other models has this feature.
\end{singlespace}
\end{footnotesize}
\end{table}
  

\begin{table}[h!]
\begin{center}
\caption{Statistical tests comparing strategic \aisub{s} vs other models ($\varepsilon = 0.05$)} \label{tab:stat_bounds_05}
\scriptsize
\input{tables/prop_ll_stats_05.tex}
\end{center}
\begin{footnotesize}
\begin{singlespace}
\vspace*{-0.05in}
\emph{Notes:} This table shows the results of the statistical tests comparing the strategic \aisub{s} to the other models for $\varepsilon = 0.05$.
No smoothing is applied to the cognitive hierarchy model.
The first column shows the comparison model.
The second presents $\bar \Lambda_S$ with bootstrap confidence intervals comparing the strategic \aisub{s} to the other models.
The third and fourth columns present p-values for the Wilcoxon signed-rank test and random-sign permutation test, respectively.
The fifth column presents the proportion of games for which the strategic \aisub{s} is the best predictor with its 95\% Clopper-Pearson interval. \starlanguage
\end{singlespace}
\end{footnotesize}
\end{table}

\begin{table}[h!]
\begin{center}
\caption{Statistical tests comparing optimized vs other models ($\varepsilon = 0.1$)} \label{tab:stat_bounds_10}
\scriptsize
\input{tables/prop_ll_stats_10.tex}
\end{center}
\begin{footnotesize}
\begin{singlespace}
\vspace*{-0.05in}
\emph{Notes:} This table shows the results of the statistical tests comparing the strategic \aisub{s} to the other models for $\varepsilon = 0.1$.
No smoothing is applied to the cognitive hierarchy model.
The first column shows the comparison model.
The second presents $\bar \Lambda_S$ with bootstrap confidence intervals comparing the strategic \aisub{s} to the other models.
The third and fourth columns present p-values for the Wilcoxon signed-rank test and random-sign permutation test, respectively.
The fifth column presents the proportion of games for which the strategic \aisub{s} is the best predictor with its 95\% Clopper-Pearson interval. \starlanguage
\end{singlespace}
\end{footnotesize}
\end{table}

\begin{table}[h!]
\begin{center}
\caption{Statistical tests comparing optimized vs other models ($\varepsilon = 0.2$)} \label{tab:stat_bounds_20}
\scriptsize
\input{tables/prop_ll_stats_20.tex}
\end{center}
\begin{footnotesize}
\begin{singlespace}
\vspace*{-0.05in}
\emph{Notes:} This table shows the results of the statistical tests comparing the strategic \aisub{s} to the other models for $\varepsilon = 0.2$.
No smoothing is applied to the cognitive hierarchy model.
The first column shows the comparison model.
The second presents $\bar \Lambda_S$ with bootstrap confidence intervals comparing the strategic \aisub{s} to the other models.
The third and fourth columns present p-values for the Wilcoxon signed-rank test and random-sign permutation test, respectively.
The fifth column presents the proportion of games for which the strategic \aisub{s} is the best predictor with its 95\% Clopper-Pearson interval. \starlanguage
\end{singlespace}
\end{footnotesize}
\end{table}

\begin{table}[h!]
\begin{center}
\caption{Statistical tests comparing optimized vs other models ($\varepsilon = 0.3$)} \label{tab:stat_bounds_30}
\scriptsize
\input{tables/prop_ll_stats_30.tex}
\end{center}
\begin{footnotesize}
\begin{singlespace}
\vspace*{-0.05in}
\emph{Notes:} This table shows the results of the statistical tests comparing the strategic \aisub{s} to the other models for $\varepsilon = 0.3$.
No smoothing is applied to the cognitive hierarchy model.
The first column shows the comparison model.
The second presents $\bar \Lambda_S$ with bootstrap confidence intervals comparing the strategic \aisub{s} to the other models.
The third and fourth columns present p-values for the Wilcoxon signed-rank test and random-sign permutation test, respectively.
The fifth column presents the proportion of games for which the strategic \aisub{s} is the best predictor with its 95\% Clopper-Pearson interval. \starlanguage
\end{singlespace}
\end{footnotesize}
\end{table}


\begin{figure}[h!]
\begin{center}
\begin{minipage}{\textwidth}
\caption{Relative predictive accuracy of the strategic \aisub{} vs other models for different game types ($\varepsilon = 0.05$)} 
\label{fig:game_level_plots_eps_005}
\vspace*{-0.1in}
\includegraphics[width=\textwidth]{plots/game_level_plots_eps_005.pdf} 
\end{minipage}
\end{center}
\begin{footnotesize}
\begin{singlespace}
\vspace*{-0.05in}
\emph{Notes:} This figure shows the proportion of games for which strategic \aisub{s} are the best predictor of initial play, separated by bonus rule (top panel) and points rule (bottom panel). 
All distributions except the cognitive hierarchy model are smoothed with $\varepsilon = 0.05$. 
The vertical dashed line represents equal performance (50-50 split). 
Green indicates that strategic \aisub{} significantly outperforms the reference model in more than 50\% of games, red indicates significantly worse performance in more than 50\% of games, and grey indicates no significant difference. 
Error bars show 95\% Clopper-Pearson confidence intervals.
\end{singlespace}
\end{footnotesize}
\end{figure}

\begin{figure}[h!]
\begin{center}
\begin{minipage}{\textwidth}
\caption{Relative predictive accuracy of the strategic \aisub{} vs other models for different game types ($\varepsilon = 0.1$)} 
\label{fig:game_level_plots_eps_01}
\vspace*{-0.1in}
\includegraphics[width=\textwidth]{plots/game_level_plots_eps_01.pdf} 
\end{minipage}
\end{center}
\begin{footnotesize}
\begin{singlespace}
\vspace*{-0.05in}
\emph{Notes:} This figure shows the proportion of games for which strategic \aisub{s} are the best predictor of initial play, separated by bonus rule (top panel) and points rule (bottom panel). 
All distributions except the cognitive hierarchy model are smoothed with $\varepsilon = 0.1$.
The vertical dashed line represents equal performance (50-50 split). 
Green indicates that strategic \aisub{} significantly outperforms the reference model in more than 50\% of games, red indicates significantly worse performance in more than 50\% of games, and grey indicates no significant difference. 
Error bars show 95\% Clopper-Pearson confidence intervals.
\end{singlespace}
\end{footnotesize}
\end{figure}

\begin{figure}[h!]
\begin{center}
\begin{minipage}{\textwidth}
\caption{Relative predictive accuracy of the strategic \aisub{} vs other models for different game types ($\varepsilon = 0.2$)} 
\label{fig:game_level_plots_eps_02}
\vspace*{-0.1in}
\includegraphics[width=\textwidth]{plots/game_level_plots_eps_02.pdf} 
\end{minipage}
\end{center}
\begin{footnotesize}
\begin{singlespace}
\vspace*{-0.05in}
\emph{Notes:} This figure shows the proportion of games for which strategic \aisub{s} are the best predictor of initial play, separated by bonus rule (top panel) and points rule (bottom panel). 
All distributions except the cognitive hierarchy model are smoothed with $\varepsilon = 0.2$.
The vertical dashed line represents equal performance (50-50 split). Green indicates that strategic \aisub{} significantly outperforms the reference model in more than 50\% of games, red indicates significantly worse performance in more than 50\% of games, and grey indicates no significant difference. 
Error bars show 95\% Clopper-Pearson confidence intervals.
\end{singlespace}
\end{footnotesize}
\end{figure}

\begin{figure}[h!]
\begin{center}
\begin{minipage}{\textwidth}
\caption{Relative predictive accuracy of the strategic \aisub{} vs other models for different game types ($\varepsilon = 0.3$)} 
\label{fig:game_level_plots_eps_03}
\vspace*{-0.1in}
\includegraphics[width=\textwidth]{plots/game_level_plots_eps_03.pdf} 
\end{minipage}
\end{center}
\begin{footnotesize}
\begin{singlespace}
\vspace*{-0.05in}
\emph{Notes:} This figure shows the proportion of games for which strategic \aisub{s} are the best predictor of initial play, separated by bonus rule (top panel) and points rule (bottom panel). 
All distributions except the cognitive hierarchy model are smoothed with $\varepsilon = 0.3$.
The vertical dashed line represents equal performance (50-50 split). Green indicates that strategic \aisub{} significantly outperforms the reference model in more than 50\% of games, red indicates significantly worse performance in more than 50\% of games, and grey indicates no significant difference. 
Error bars show 95\% Clopper-Pearson confidence intervals.
\end{singlespace}
\end{footnotesize}
\end{figure}

\newpage \clearpage

\begin{table}[h!] 
\begin{center}
\begin{minipage}{1 \linewidth}                          
\caption{Log-likelihood ratio regressions across game types ($\varepsilon = 0.05$)}  
\label{tab:game_level_eps_005}     
\centering                                                  
\scriptsize
\renewcommand{\arraystretch}{1.2}
\input{tables/game_level_eps_005.tex}
\end{minipage}
\end{center}
\begin{footnotesize}
\begin{singlespace}
\vspace*{-0.15in}
\emph{Notes:} Each column reports regression estimates where the dependent variable is the log-likelihood ratio of the strategic \aisub{} to the other models.
The independent variables are the game features from Table~\ref{tab:game_parameters} (excluding the bonus rule and points rule).
We report heteroskedasticity-robust standard errors.
\starlanguage{}
\end{singlespace}
\end{footnotesize}
\end{table}

\begin{table}[h!] 
\begin{center}
\begin{minipage}{1 \linewidth}                          
\caption{Log-likelihood ratio regressions across game types ($\varepsilon = 0.1$)}  
\label{tab:game_level_eps_01}     
\centering                                                
\scriptsize
\renewcommand{\arraystretch}{1.2}
\input{tables/game_level_eps_01.tex}
\end{minipage}
\end{center}
\begin{footnotesize}
\begin{singlespace}
\vspace*{-0.15in}
\emph{Notes:} Each column reports regression estimates where the dependent variable is the log-likelihood ratio of the strategic \aisub{} to the other models.
The independent variables are the game features from Table~\ref{tab:game_parameters} (excluding the bonus rule and points rule).
We report heteroskedasticity-robust standard errors.
\starlanguage{}
\end{singlespace}
\end{footnotesize}
\end{table}

\begin{table}[h!] 
\begin{center}
\begin{minipage}{1 \linewidth}                          
\caption{Log-likelihood ratio regressions across game types ($\varepsilon = 0.2$)}  
\label{tab:game_level_eps_02}     
\centering                                            
\scriptsize
\renewcommand{\arraystretch}{1.2}
\input{tables/game_level_eps_02.tex}
\end{minipage}
\end{center}
\begin{footnotesize}
\begin{singlespace}
\vspace*{-0.15in}
\emph{Notes:} Each column reports regression estimates where the dependent variable is the log-likelihood ratio of the strategic \aisub{} to the other models.
The independent variables are the game features from Table~\ref{tab:game_parameters} (excluding the bonus rule and points rule).
We report heteroskedasticity-robust standard errors.
\starlanguage{}
\end{singlespace}
\end{footnotesize}
\end{table}

\begin{table}[h!] 
\begin{center}
\begin{minipage}{1 \linewidth}                          
\caption{Log-likelihood ratio regressions across game types ($\varepsilon = 0.3$)}  
\label{tab:game_level_eps_03}     
\centering                                           
\scriptsize
\renewcommand{\arraystretch}{1.2}
\input{tables/game_level_eps_03.tex}
\end{minipage}
\end{center}
\begin{footnotesize}
\begin{singlespace}
\vspace*{-0.15in}
\emph{Notes:} Each column reports regression estimates where the dependent variable is the log-likelihood ratio of the strategic \aisub{} to the other models.
The independent variables are the game features from Table~\ref{tab:game_parameters} (excluding the bonus rule and points rule).
We report heteroskedasticity-robust standard errors.
\starlanguage{}
\end{singlespace}
\end{footnotesize}
\end{table}

